\documentclass{article}
\usepackage[UTF8]{ctex}
\usepackage{amsmath,amssymb,amsfonts}
\usepackage{enumitem}
\usepackage{geometry}
\geometry{a4paper, margin=1in}
% === 错题本排版优化（不参与题目标识） ===
\setlength{\parindent}{0pt}
\setlength{\parskip}{0.35em plus 0.10em minus 0.08em}
\setlength{\emergencystretch}{3em}
\setlist{itemsep=0.25em, topsep=0.40em, parsep=0pt}
\allowdisplaybreaks[2]
% === 错题本排版优化结束 ===
\begin{document}
\section*{东北师范大学2024年数学分析真题}
\section*{78．东北师范大学2024年数学分析试卷}

\subsection{题目1}
一．计算题．每题6分，共30分．

\begin{itemize}
  \item[（1）] \(\lim _{n \rightarrow \infty}\left(\frac{1}{n}+\frac{1}{n+1}+\cdots+\frac{1}{2 n+1}\right)\) ．
  \item[（2）] \(\lim _{x \rightarrow+\infty} \sqrt{x} \int_{x}^{x+1} \frac{\mathrm{~d} t}{\sqrt{t+\ln t}}\) ．
  \item[（3）] \(\lim _{n \rightarrow \infty}\left(\frac{1}{2}+\frac{3}{2^{2}}+\frac{5}{2^{3}}+\cdots+\frac{2 n-1}{2^{n}}\right)\) ．
  \item[（4）] \(\int_{0}^{1} \mathrm{~d} x \int_{0}^{1}\left|x y-\frac{1}{4}\right| \mathrm{d} y\) ．
  \item[（5）] 求球面 \(x=\sin \varphi \cos \theta, y=\sin \varphi \sin \theta, z=\cos \varphi\) 在对应 \(\theta=\varphi=\frac{\pi}{4}\) 处的切平面和法线方程．
\end{itemize}

\subsection{题目2}
二．（15 分）由方程组

\[
\left\{\begin{array}{l}
x^{2}+y^{2}-u v=0 \\
x y+u^{2}-v^{2}=0
\end{array}\right.
\]

能否确定 \(u, v\) 为 \(x\) 与 \(y\) 的函数？在能确定隐函数的条件下，求 \(u_{x}, v_{x}, u_{y}, v_{y}\) ．

\subsection{题目3}
三．（15 分）求函数 \(f(x, y)=\sin x+\sin y+\sin (x+y)\) 在 \((x, y) \in[0,2 \pi] \times[0,2 \pi]\) 的所有极值点．

\subsection{题目4}
四．（15 分）计算曲面积分

\[
I=\iint_{S}\left(z^{2}+x\right) \mathrm{d} y \mathrm{~d} z+\sqrt{z} \mathrm{~d} x \mathrm{~d} y
\]

其中 \(S\) 为抛物面 \(z=\frac{1}{2}\left(x^{2}+y^{2}\right)\) 在 \(z=0\) 和 \(z=\pi\) 之间的部分，定向取下侧．

\subsection{题目5}
五．（15 分）设函数 \(f(x)\) 在 \((-\infty,+\infty)\) 上有任意阶导数，且对任意实数 \(x\) 及 \(n=0,1,2, \cdots\) ，满足

\[
\left|f^{(n)}(x)\right| \leq n!|x| .
\]

证明：\(f(x)=0, x \in(-\infty,+\infty)\) ．

\subsection{题目6}
六。设 \(D\) 为平面区域，\(u(x, y) \in C^{2}(D)\) ，证明：\(u(x, y)\) 为调和函数，即 \(u\) 满足 \(\frac{\partial^{2} u}{\partial x^{2}}+\frac{\partial^{2} u}{\partial y^{2}}=0\) 的充要条件是：对 \(D\) 内任一圆周 \(L\) ，且 \(L\) 所围圆属于 \(D\) ，都有 \(\int_{L} \frac{\partial u}{\partial n} \mathrm{~d} s=0\) 。

\subsection{题目7}
七．（15 分）设 \(\left\{a_{n}\right\}\) 是正的单调递增序列，证明：级数 \(\sum_{n=1}^{\infty}\left(1-\frac{a_{n}}{a_{n+1}}\right)\) 当 \(\left\{a_{n}\right\}\) 有界时收敛，当 \(\left\{a_{n}\right\}\)无界时发散。

\begin{itemize}
  \item[八．] （15 分）令 \(f(t)=\int_{1}^{+\infty} \frac{\cos x t}{1+x^{2}} \mathrm{~d} x\) ．证明：
  \item[（1）] 所定义的函数 \(f(t)\) 在 \(\mathbb{R}\) 上一致连续．
  \item[（2）] \(f(t) \in C(-\infty,+\infty)\) ．
  \item[（3）] \(\lim _{t \rightarrow+\infty} f(t)=0\) ．
\end{itemize}

\subsection{题目8}
九．（15 分）设 \(\left\{f_{n}\right\}\) 是 \([a, b]\) 上连续的函数列，对任意的 \(x \in[a, b]\) ，有 \(f_{n}(x) \leq f_{n+1}(x), n=1,2, \cdots\) ．证明：如果 \(\left\{f_{n}\right\}\) 在 \([a, b]\) 上收敛于连续函数 \(f\) ，那么 \(\left\{f_{n}\right\}\) 必定一致收敛于 \(f\) 。
\end{document}
